% year: תשפ"ד
% type: מבחן
% semester: א
% moed: ב
% number: 7
% points: 14
% categories: improper-integrals, func-limits
% source: exams/מבחנים/תשפד/מועד ב תשפד.pdf

\begin{question}
בדקו התכנסות של האיטגרל הלא-אמיתי הבא:
\[ \int_2^\infty \frac{\sqrt[5]{x}}{3x^2+x+100}\,dx \]
\end{question}

\begin{hint}
לאינסוף, הפונקציה מתנהגת כמו $\frac{x^{1/5}}{3x^2}$. השוו (במבחן ההשוואה הגבולי) ל-$\frac{1}{x^{9/5}}$.
\end{hint}

\begin{solution}
הפונקציה $f(x)=\frac{\sqrt[5]{x}}{3x^2+x+100}$ רציפה וחיובית על $[2,\infty)$ (המכנה חיובי), ולכן האינטגרל לא-אמיתי רק בגלל הגבול העליון, ומותר להשתמש במבחני ההשוואה לפונקציות אי-שליליות.

נשווה ל-$g(x)=\frac{1}{x^{9/5}}=\frac{x^{1/5}}{x^2}$, שגם היא חיובית על $[2,\infty)$:
\[ \lim_{x\to\infty}\frac{f(x)}{g(x)}=\lim_{x\to\infty}\frac{x^{1/5}\cdot x^{2}}{x^{1/5}(3x^2+x+100)}=\lim_{x\to\infty}\frac{1}{3+\frac1x+\frac{100}{x^2}}=\frac13. \]
הגבול סופי וחיובי, ולכן לפי \textbf{מבחן ההשוואה הגבולי} האינטגרלים $\int_2^\infty f$ ו-$\int_2^\infty g$ מתכנסים או מתבדרים יחד.

האינטגרל $\int_2^\infty\frac{dx}{x^p}$ מתכנס אם ורק אם $p>1$; כאן $p=\frac95>1$, ובאופן מפורש
\[ \int_2^\infty x^{-9/5}dx=\lim_{R\to\infty}\left[-\frac54x^{-4/5}\right]_2^R=\frac54\cdot2^{-4/5}<\infty. \]

\textbf{תשובה:} האינטגרל $\int_2^\infty\frac{\sqrt[5]{x}}{3x^2+x+100}\,dx$ \textbf{מתכנס}.
\end{solution}
